\documentclass[11pt]{article}

\usepackage[margin=1in]{geometry}
\usepackage{amsmath,amssymb,amsthm,mathtools}
\usepackage{mathrsfs}
\usepackage{bm}
\usepackage[T1]{fontenc}
\usepackage{lmodern}
\usepackage{microtype}
\usepackage{booktabs}
\usepackage{enumitem}
\usepackage{hyperref}

\hypersetup{colorlinks=true,linkcolor=blue,citecolor=blue,urlcolor=blue}

\newtheorem{theorem}{Theorem}[section]
\newtheorem{proposition}[theorem]{Proposition}
\newtheorem{corollary}[theorem]{Corollary}
\newtheorem{lemma}[theorem]{Lemma}
\theoremstyle{definition}
\newtheorem{definition}[theorem]{Definition}
\newtheorem{remark}[theorem]{Remark}

\newcommand{\kk}{\Bbbk}
\newcommand{\im}{\operatorname{im}}
\newcommand{\kerp}{\operatorname{ker}}
\newcommand{\End}{\operatorname{End}}
\newcommand{\Hom}{\operatorname{Hom}}
\newcommand{\rankp}{\operatorname{rank}}
\newcommand{\id}{\operatorname{id}}
\newcommand{\Spec}{\operatorname{Spec}}
\newcommand{\Aeff}{\mathscr A_{\mathrm{eq}}^{\mathrm{eff}}}
\newcommand{\Aeq}{\mathscr A_{\mathrm{eq}}}
\newcommand{\Neq}{N_{\mathrm{eq}}}
\newcommand{\Geq}{\Gamma_{\mathrm{eq}}}
\newcommand{\Gcan}{\mathfrak G_{\mathrm{can}}}
\newcommand{\Csyz}{\mathscr C_{\mathrm{syz}}}
\newcommand{\forced}{\mathrm{forced}}

\title{General Spectral Floors from Marked Deformation Presentations\\
\large Equation-Operator Reconstruction, Syzygy Propagation, and Normalized-Tilt Invariants}
\author{Research closeout draft, Paper III}
\date{September 5, 2026}

\begin{document}
\maketitle

\begin{abstract}
We give a finite algebraic pipeline for extracting spectral information that survives a prescribed
class of normalized tilts. The input is a marked deformation presentation consisting of a finite
filtered relation resolution, a canonical family of chain operators, and a relation-seed map. The
canonical operators first generate an effective equation-operator algebra on zeroth homology. This
algebra generates the smallest invariant normalized source, which is combined with source and
landing maps to form a forced sector. A general orbit theorem then identifies the characteristic
and minimal polynomial divisors common to all normalized tilts with the corresponding polynomials
on the forced sector. We also give the syzygy-to-propagator bridge, finite closure bounds, and exact
finite models. The result is deliberately scoped: the operators are canonical relative to the
marked presentation, and the spectral floor is a forced divisor rather than a full unmarked
spectrum. Filtration alone cannot manufacture the distinguished operator family.
\end{abstract}

\tableofcontents

\section{Introduction and scope}

The purpose of this paper is to close a finite spectral-floor program for inverse-Leibniz type
deformation data. The central question is not to recover every eigenvalue of one chosen operator,
but to determine which polynomial factors cannot be removed when the operator is allowed to vary
through a controlled normalized-tilt orbit.

The paper collects five successive reductions. Version v0.57 isolates the normalized-tilt floor.
Version v0.58 reconstructs the forced sector from landing and normalization data. Version v0.59
generates the normalized source from relation seeds rather than taking it as an external input.
Version v0.60 explains how source-side propagators descend from an equation-operator action on a
syzygy presentation. Version v0.61 removes the separately declared equation algebra and reconstructs
the effective algebra from canonical chain operators.

The resulting pipeline is
\begin{equation}
\boxed{
(P_\bullet,\Gcan,R_0)
\longmapsto \Aeff
\longmapsto \Neq
\longmapsto \Geq
\longmapsto (K,A|_K)
\longmapsto (\chi_{\forced},\mu_{\forced})
}.
\label{eq:pipeline}
\end{equation}
Here \(P_\bullet\) is a finite filtered relation resolution, \(\Gcan\) is a specified family of
canonical chain operators, and \(R_0\) is a specified relation-seed map. The adjective ``canonical''
is relative to this marked presentation. It does not mean that a bare filtration or a bare chain
complex determines the operators.

Our conclusions are algebraic and finite. Exact matrix computations illustrate the statements, but
do not replace the hypotheses of the theorems or establish an infinite-dimensional realization.
Likewise, the forced divisor is not the full spectrum of an arbitrary representative. The paper
stops at the marked deformation category and does not claim an ordinary unmarked
quasi-isomorphism invariant.

\section{Normalized tilts and spectral floors}

Let \(\kk\) be an infinite field, let \(R\) and \(F\) be finite-dimensional \(\kk\)-vector spaces, and let
\[
D:R\longrightarrow F, \qquad N\subseteq F, \qquad q:F\longrightarrow F/N
\]
be linear data. Define the normalized kernel
\begin{equation}
K:=D^{-1}(N),
\qquad A\in\End(R).
\label{eq:K}
\end{equation}
The allowed normalized tilts are parametrized by \(\tau\in\Hom(\im(qD),R)\), with
\begin{equation}
A_\tau:=A+\tau qD.
\label{eq:tilt}
\end{equation}
The perturbation is invisible on \(K\), so every representative has the same restriction to
\(K\).

\begin{lemma}[Exact variation space]
The set of all perturbations obtained from normalized tilts is
\begin{equation}
\{\tau qD:\tau\in\Hom(\im(qD),R)\}
=\{T\in\End(R):T|_K=0\}.
\label{eq:variation}
\end{equation}
\end{lemma}

\begin{proof}
Every \(\tau qD\) vanishes on \(K\). Conversely, if \(T|_K=0\), define \(\tau(qD r):=Tr\).
If \(qDr=qDr'\), then \(r-r'\in K\), and hence \(T(r-r')=0\). Thus \(\tau\) is well-defined
and \(T=\tau qD\).
\end{proof}

\begin{theorem}[Normalized-tilt spectral floor]
Assume \(A(K)\subseteq K\). If \(K\ne 0\), then for every normalized tilt \(A_\tau\), the characteristic polynomial of
\(A|_K\) divides \(\chi_{A_\tau}\), and the minimal polynomial of \(A|_K\) divides
\(\mu_{A_\tau}\). Moreover,
\begin{align}
\gcd_{\tau}\chi_{A_\tau}&=\chi_{A|_K},
\\
\gcd_{\tau}\mu_{A_\tau}&=\mu_{A|_K},
\end{align}
where the gcd is taken monically over all allowed tilts.
If \(K=0\), the characteristic floor is the unit polynomial and the minimal floor is also the
unit polynomial by convention.
\label{thm:floor}
\end{theorem}

\begin{proof}
Choose a complement \(X\) with \(R=K\oplus X\). Since every perturbation vanishes on \(K\),
the matrix of \(A_\tau\) has block form
\[
\begin{pmatrix}
A|_K & *\\
0 & B_\tau
\end{pmatrix}.
\]
Consequently \(\chi_{A_\tau}=\chi_{A|_K}\chi_{B_\tau}\), and the restriction polynomial
\(\chi_{A|_K}\) divides every characteristic polynomial. The same block form gives
\(\mu_{A|_K}\mid\mu_{A_\tau}\).

By the exact variation lemma, every endomorphism of \(X\) can be chosen as the lower-right block
of a perturbation, and the off-diagonal block can be chosen arbitrarily as well. In particular,
choose a tilt with zero off-diagonal block and with \(B_\tau=cI_X\), where \(c\) is outside the
finite set of eigenvalues of \(A|_K\). Then the complementary characteristic factor is coprime
to \(\chi_{A|_K}\), so no nonconstant factor beyond \(\chi_{A|_K}\) survives the gcd. For minimal
polynomials, choose a scalar complementary block with eigenvalue avoiding the roots of
\(\mu_{A|_K}\). The lower bound is already supplied by restriction, and the two bounds coincide.
\end{proof}

\begin{corollary}[Jordan-depth interpretation]
If
\[
\chi_{A|_K}(S)=\prod_\lambda(S-\lambda)^{c_\lambda},
\qquad
\mu_{A|_K}(S)=\prod_\lambda(S-\lambda)^{j_\lambda},
\]
then \(c_\lambda\) and \(j_\lambda\) are respectively the algebraic-multiplicity and
Jordan-depth floors that survive every normalized tilt.
\end{corollary}

\begin{remark}
The proof uses the full normalized variation space in \eqref{eq:variation}. If an application
restricts the tilt parameters further, the gcd can be larger. The theorem therefore identifies the
floor for the stated orbit, not for every conceivable deformation.
\end{remark}

\section{Forced-sector extraction from finite landing data}

We now remove \(K\) as an external input. Let \(M\) be a finite-dimensional state space, let
\(T\in\End(M)\), and let \(\Lambda:M\to R\) be a landing map. Retain \(D:R\to F\), a normalized
subspace \(N\subseteq F\), and \(q:F\to F/N\). Define
\begin{equation}
\Gamma:=qD\Lambda:M\longrightarrow F/N.
\label{eq:Gamma}
\end{equation}

\begin{proposition}[Static forced sector]
The part of the normalized kernel visible through the landing map is
\[
K\cap\im\Lambda=\Lambda(\ker\Gamma),
\]
and there is a natural vector-space isomorphism
\begin{equation}
K\cap\im\Lambda\simeq\ker\Gamma/\ker\Lambda.
\label{eq:quotient}
\end{equation}
In particular,
\begin{equation}
\dim(K\cap\im\Lambda)=\rankp\Lambda-\rankp\Gamma.
\label{eq:dimension}
\end{equation}
\end{proposition}

\begin{proof}
By definition, \(\Lambda m\in K\) exactly when \(qD\Lambda m=\Gamma m=0\). Thus the image is
\(\Lambda(\ker\Gamma)\). The first isomorphism theorem gives the quotient description, and
\(\dim\ker\Gamma-\dim\ker\Lambda=\rankp\Lambda-\rankp\Gamma\).
\end{proof}

\begin{proposition}[Closure tests]
The landing quotient \(M/\ker\Lambda\) is stable under the induced recurrent action precisely when
\(\ker\Lambda\) is \(T\)-invariant. Equivalently,
\[
\rankp\begin{bmatrix}\Lambda\\\Lambda T\end{bmatrix}=\rankp\Lambda.
\]
The normalized quotient \(\ker\Gamma/\ker\Lambda\) is stable precisely when \(\ker\Gamma\) is
\(T\)-invariant, equivalently
\[
\rankp\begin{bmatrix}\Gamma\\\Gamma T\end{bmatrix}=\rankp\Gamma.
\]
When both conditions hold, and the landing map intertwines the recurrent action with the operator on
the landed space, the action induced by \(T\) on \(\ker\Gamma/\ker\Lambda\) is conjugate to the
action on \(K\cap\im\Lambda\).
\end{proposition}

\begin{proof}
The stacked-rank criterion is the standard finite-dimensional test for
\(\ker L\subseteq\ker(LT)\). Thus each displayed condition is exactly the condition that the
corresponding quotient action is well-defined. Under the stated intertwining condition, the induced
action is transported through the quotient isomorphism \eqref{eq:quotient}.
\end{proof}

If closure fails, define the stable normalized core
\begin{equation}
M_\infty:=\bigcap_{n\geq 0}\ker(\Gamma T^n).
\label{eq:stablecore}
\end{equation}
Finite dimensionality makes this intersection stationary after finitely many steps. The same
argument gives \(K_\infty\simeq M_\infty/\ker\Lambda\) whenever \(\ker\Lambda\subseteq M_\infty\), with
\[
\dim K_\infty=\rankp\Lambda-\rankp\mathcal O_N,
\qquad
\mathcal O_N:=\begin{bmatrix}\Gamma\\\Gamma T\\\Gamma T^2\\\vdots\end{bmatrix}.
\]
This is the stable version used when the landing data are not closed at one step.

\begin{corollary}[Cyclic PID form]
Suppose \(M=R/(c)\) is cyclic over a PID and
\[
\ker\Lambda=d_\Lambda R/(c),
\qquad
\ker\Gamma=d_\Gamma R/(c).
\]
Then \(d_\Gamma\mid d_\Lambda\),
\[
K\cap\im\Lambda\simeq R/(d_\Lambda/d_\Gamma),
\]
and \(d_\Lambda/d_\Gamma\) is the cyclic forced polynomial.
\end{corollary}

\section{Generating the normalized source}

The normalized source should not be postulated when it can be generated from relation data. Let
\(W_{\rm rel}\) be a finite relation-seed space, let
\[
R_0:W_{\rm rel}\longrightarrow F
\]
be the seed map, and let \(C_1,\ldots,C_s\in\End(F)\) be source-side constraint propagators.
Starting from \(N^{(0)}=\im R_0\), define
\begin{equation}
N^{(r+1)}=N^{(r)}+\sum_{i=1}^s C_iN^{(r)},
\qquad
\Neq:=\bigcup_{r\geq 0}N^{(r)}.
\label{eq:closure}
\end{equation}
Because \(F\) is finite-dimensional, the increasing sequence stabilizes.

\begin{proposition}[Minimal invariant source]
The stable space in \eqref{eq:closure} is
\begin{equation}
\Neq=\kk\langle C_1,\ldots,C_s\rangle\im R_0,
\label{eq:Neq}
\end{equation}
and it is the unique smallest subspace of \(F\) that contains \(\im R_0\) and is invariant under
every \(C_i\).
\end{proposition}

\begin{proof}
Induction gives \(N^{(r)}\) equal to the span of all words in the \(C_i\) of length at most
\(r\) applied to \(\im R_0\). The stable union is therefore the right-hand side of
\eqref{eq:Neq}. Any invariant subspace containing the seed contains every such word, proving
minimality.
\end{proof}

\begin{proposition}[Functoriality]
Suppose \(\Phi_F:F\to F'\) and \(\Phi_W:W_{\rm rel}\to W'_{\rm rel}\) satisfy
\[
\Phi_FR_0=R_0'\Phi_W,
\qquad
\Phi_FC_i=C_i'\Phi_F.
\]
Then \(\Phi_F(\Neq)\subseteq N_{\rm eq}'\). If \(\Phi_F\) is an isomorphism intertwining all
data, it carries \(\Neq\) isomorphically to \(N_{\rm eq}'\).
\end{proposition}

The forced-sector construction is now applied with \(N=\Neq\), giving
\[
\Geq=q_{\rm eq}D\Lambda,
\qquad
K\simeq\ker\Geq/\ker\Lambda,
\]
subject to the closure or stable-core hypotheses of Section 3. In this form no independently
chosen normalized source remains in the pipeline.

\section{Syzygies and the constraint propagator}

The source propagators themselves must come from an equation action. Let \(\Aeq\) temporarily denote
an algebra acting on a finite relation module \(F_{\rm rel}\). Let
\[
P_1\xrightarrow{d_1}P_0\xrightarrow{\varepsilon}F_{\rm rel}\longrightarrow 0
\]
be an exact presentation in left \(\Aeq\)-modules. Equivalently,
\(F_{\rm rel}=H_0(P_\bullet)\) for the two-term complex.

\begin{proposition}[Syzygy descent]
If \(a\in\Aeq\) acts on the resolution by chain maps, then it descends to an endomorphism
\(\bar\rho(a)\in\End(F_{\rm rel})\). Via \(\varepsilon\), the induced source constraint
propagator is
\begin{equation}
C(a)=\bar\varepsilon\,\bar\rho(a)\,\bar\varepsilon^{-1}
\quad\text{on the identified relation image},
\label{eq:Csyzygy}
\end{equation}
and the effective constraint algebra is
\begin{equation}
\Csyz=\im\bigl(C:\Aeq\to\End(F_{\rm rel})\bigr)
\simeq\Aeq/\operatorname{Ann}(F_{\rm rel}).
\label{eq:Csyz}
\end{equation}
\end{proposition}

\begin{proof}
The chain-map condition preserves \(\im d_1\), so the action on \(P_0\) descends to
\(P_0/\im d_1=H_0(P_\bullet)\). Transport along the presentation isomorphism gives
\eqref{eq:Csyzygy}. The kernel of the resulting representation is exactly the annihilator of the
relation module, which proves \eqref{eq:Csyz}.
\end{proof}

Changing generators or replacing the presentation by an isomorphic resolution changes matrices but
not the induced representation on \(H_0\). An ambient extension of \(C(a)\) outside the relation
image is noncanonical and has no role in \(\Neq\). This is the precise syzygy-to-propagator bridge
needed by the source-closure step.

\section{Reconstructing the effective equation algebra}

The temporary algebra \(\Aeq\) in Section 5 is now removed from the external input. Let
\(P_\bullet\) be a finite filtered relation resolution and let
\[
\Gcan=\{G_1,\ldots,G_t\}
\]
be the specified canonical chain-operator family. Assume each operator is a chain map:
\[
dG_j=G_jd.
\]
Write \(H_0=H_0(P_\bullet)\). Each \(G_j\) descends to \(\bar G_j\in\End(H_0)\).

\begin{definition}[Effective equation-operator algebra]
The effective algebra of the marked presentation is
\begin{equation}
\Aeff:=\kk\langle\bar G:G\in\Gcan\rangle
\subseteq\End(H_0).
\label{eq:Aeff}
\end{equation}
Equivalently, if \(\kk\langle X_G:G\in\Gcan\rangle\) is the free unital algebra, evaluation gives
\begin{equation}
\im(\operatorname{ev}_P)
\simeq
\kk\langle X_G:G\in\Gcan\rangle/\ker(\operatorname{ev}_P).
\label{eq:freeeval}
\end{equation}
\end{definition}

\begin{proposition}[Minimality and naturality]
The algebra \(\Aeff\) is the smallest unital subalgebra of \(\End(H_0)\) containing every descended
canonical operator. If two marked presentations are related by a filtered chain comparison \(\Phi\)
and homotopies \(H_G\) satisfying
\[
\Phi G-G'\Phi=dH_G+H_Gd,
\]
then \(\Phi\) induces a map on \(H_0\) that intertwines the two effective algebras. In particular,
the effective algebras are conjugate when the induced map is an isomorphism.
\end{proposition}

\begin{proof}
The first claim is the definition of the generated unital algebra. The homotopy equation becomes
\(\bar\Phi\bar G=\bar G'\bar\Phi\) after passing to homology, since boundaries vanish. Thus
\(\bar\Phi\) carries every word in the first operator family to the corresponding word in the
second family, and the same holds for linear combinations and products.
\end{proof}

Raw actions on contractible summands may differ between resolutions. They do not affect \(H_0\),
and hence they do not affect \(\Aeff\). The collapse of the preceding section is therefore
\[
\Csyz=\Aeff,
\qquad
\Neq=\Aeff\,\im R_0,
\]
with the representation understood on the relevant relation module.

\begin{remark}[Sharp no-go]
The data of a filtration and a bare chain complex do not manufacture a distinguished family
\(\Gcan\). The effective-algebra theorem is consequently a theorem about marked deformation
presentations. Removing the marks is a different problem and is not part of this paper.
\end{remark}

\section{Finite closure and filtered triangularity}

All closures in this paper are finite-dimensional. If \(H_0\) has dimension \(h\), the algebra
generated by a finite operator family lies in \(\End(H_0)\), of dimension \(h^2\). A naive closure
algorithm that adjoins products and retains their span can therefore have at most \(h^2\) strict
rank increases. The same finite argument gives termination for source closure in \(F\).

For a finite descending filtration such that positive-filtration operators strictly raise the
filtration degree, those operators form a nilpotent ideal after passing to the associated finite
filtered representation.
This triangularity explains why the effective algebra can be computed by finite rank tests. It does
not turn a finite computation into an infinite-dimensional spectral theorem.

The complete finite procedure is:
\begin{enumerate}[leftmargin=2em]
\item compute \(H_0(P_\bullet)\) and descend every \(G\in\Gcan\);
\item close the descended family under products to obtain \(\Aeff\);
\item close \(\im R_0\) under \(\Aeff\) to obtain \(\Neq\);
\item form \(\Geq=q_{\rm eq}D\Lambda\) and test the two kernel-invariance conditions;
\item compute the induced operator on \(K\simeq\ker\Geq/\ker\Lambda\);
\item report its characteristic and minimal polynomials as the forced spectral floors.
\end{enumerate}

\section{Exact finite models and verification}

\subsection{Equation-operator reconstruction model}

Take \(H_0=\kk^4\), let \(S=J_4(0)\) be the nilpotent shift, and let
\[
W=\operatorname{diag}(0,1,2,3).
\]
Then
\[
S^4=0,
\qquad
[W,S]=S.
\]
The algebra generated by \(S\) and \(W\) is the lower-triangular matrix algebra in this basis,
with dimension \(10\). The shift-only algebra \(\kk[S]\) has dimension \(4\). A change of basis
preserves these dimensions and relations after conjugation. Adding a contractible summand changes
raw chain matrices but leaves the induced \(H_0\) action unchanged. These checks verify precisely
the claims that matter for v0.61; they do not assert that arbitrary chain complexes carry such
operators canonically.

\subsection{Earlier stages of the pipeline}

The exact finite checks recorded in the Part III folder give the following values:
\begin{center}
\begin{tabular}{@{}lll@{}}
\toprule
Version & finite output & role \\
\midrule
v0.57 & \(\chi_{\forced}=\mu_{\forced}=(S+2)^2(S-3)\), tilt dimension \(18\) & floor theorem \\
v0.58 & forced-sector dimension \(3\); cyclic floor \((S+2)^2\) & quotient extraction \\
v0.59 & ranks \(2,4,5,5\); \(\dim\Neq=5\) & source closure \\
v0.60 & syzygy rank \(4\); quotient dimension \(4\); \(\lambda^4\) & propagator bridge \\
v0.61 & \(\dim\Aeff=10\), shift-only dimension \(4\) & operator reconstruction \\
\bottomrule
\end{tabular}
\end{center}

The five verifier scripts are finite symbolic or exact matrix checks. Their PASS outputs are evidence
for the listed models and identities. They are not substitutes for the hypotheses in
Theorems~\ref{thm:floor} and the propositions above.

\section{No-go boundary and final scope}

The program has a sharp boundary. A bare filtered chain complex determines its homology and its
filtration data, but it does not select a distinguished endomorphism family. Two systems with the
same underlying finite quotient can therefore carry different operator actions, for example a
nilpotent Jordan block or a diagonal operator. Syzygies alone cannot choose between them.

The correct input is consequently the marked tuple
\[
(P_\bullet,\Gcan,R_0),
\]
where the marks are part of the deformation presentation. Within this category, the operator
algebra, normalized source, forced sector, and forced spectral floor are defined and functorial under
the stated comparisons. Outside it, this paper makes no claim of an ordinary unmarked
quasi-isomorphism invariant.

In particular, the following are deliberately future directions rather than conclusions here:
\begin{itemize}[leftmargin=2em]
\item a Natural-Operator Characterization Theorem inside a specified marked category;
\item an unmarked formal-moduli reconstruction;
\item arbitrary infinite-dimensional or full-spectrum extensions;
\item a new paper based only on renamed or overlapping constructions without a new independent
question and a new main theorem.
\end{itemize}

\section{Conclusion}

Versions v0.57--v0.61 form a complete finite general spectral-floor pipeline. Each stage removes one
external declaration: first the complementary spectrum, then the forced sector, then the normalized
source, then the source propagators, and finally the equation-operator algebra. The remaining data
are exactly the marked deformation presentation and the relation seeds. The resulting theorem is
therefore complete at the level claimed: it identifies the spectral floor forced by normalized tilts
for finite marked presentations, with exact finite verification and an explicit no-go boundary.

The paper stops here. Any continuation must introduce a genuinely new question and prove a theorem
that is not already contained in this pipeline or ruled out by the unmarked no-go boundary.

\begin{thebibliography}{9}
\bibitem{paperII}
Research closeout, \emph{Homotopy, Tilt, and Naturality in the Inverse-Leibniz Cubic Case},
Paper II, v0.01, 2026.

\bibitem{partIII}
Part III verification archive, versions v0.57--v0.61, finite exact models and verifier scripts,
2026.
\end{thebibliography}

\end{document}
